\section{Chapitre~\ref{sec:gaba}. \nt{GLUT} et \nt{GABA} ensemble}

Les énoncés logiques et dynamiques du chapitre y sont déduits de l'analyse linéaire et de la loi d'Ohm; l'expérience montre que les circuits sur lesquels ils reposent (veto retardé, inhibition directe à la suite de l'excitation, stabilisation par l'inhibition, rythme de la boucle \nt{GLUT}--\nt{GABA}) fonctionnent réellement dans le cerveau.

{\footnotesize
\setlength{\tabcolsep}{3pt}\renewcommand{\arraystretch}{1.15}
\begin{longtable}{>{\raggedright}p{0.5cm}>{\raggedright}p{4.0cm}>{\raggedright}p{5.6cm}>{\raggedright}p{2.3cm}>{\raggedright\arraybackslash}p{1.7cm}}
\toprule
N\textsuperscript{o} & Affirmation & Expérience et ce qui est mesuré & Source & Statut \\
\midrule\endfirsthead
\toprule
N\textsuperscript{o} & Affirmation & Expérience et ce qui est mesuré & Source & Statut \\
\midrule\endhead
1 & Les neurones \nt{GABA} représentent d'un cinquième à un quart des neurones du cortex & Immunomarquage du GABA dans $10$ aires du cortex du singe: environ $25\,\%$ des neurones dans $8$ aires, moins de $20\,\%$ dans le cortex visuel primaire. Revue de la diversité des neurones inhibiteurs du cortex & \cite{Hendry1987,Markram2004} & mesuré \\
2 & Ce que fait l'inhibition dépend en grande partie du site d'arrivée: dendrite, corps, lieu de naissance de l'impulsion, terminaison du fil d'un autre neurone & Revue: les classes de neurones \nt{GABA} du cortex se distinguent par leur cible sur le destinataire. Enregistrement simultané du corps et d'une dendrite d'un neurone pyramidal: l'inhibition directe est bien plus forte sur le corps que sur les dendrites. L'inhibition sur la terminaison du fil d'un autre neurone, dans la moelle épinière, réduit la libération de neurotransmetteur d'une seule ligne d'entrée (revue de mesures) & \cite{Markram2004,Pouille2001,Rudomin1999} & mesuré \\
3 & Le changement de signe par un intermédiaire coûte un retard, environ une milliseconde; l'inhibition qui suit l'excitation rétrécit la fenêtre de coïncidence & Neurones pyramidaux de l'hippocampe du rat en tranches: l'inhibition par un intermédiaire suit l'excitation directe avec un court retard, et la fenêtre dans laquelle deux entrées excitatrices se somment jusqu'au seuil est inférieure à $2\ms$. Sur les dendrites, où cette inhibition est plus faible, la fenêtre est plus large & \cite{Pouille2001} & mesuré \\
4 & Le veto $a\wedge\bar b$~\eqref{eq:veto}, avec le OU et une constante égale à un, est fonctionnellement complet (proposition~\ref{prop:gabacomplete}) & Démonstration dans le chapitre. Résultat classique: un réseau d'éléments à seuil dotés d'entrées inhibitrices réalise n'importe quelle expression logique. Énoncé sur le modèle, pas d'expérience & \cite{McCullochPitts1943} & théorie \\
5 & Le veto retardé donne un détecteur de direction du mouvement de vitesse optimale $v^*=\Delta x/d$~\eqref{eq:vstar} & Rétine du lapin: la sélectivité à la direction est créée par une inhibition qui, pour un mouvement dans la direction \og nulle\fg, éteint l'excitation. L'enregistrement séparé des entrées excitatrice et inhibitrice des cellules sélectives a montré que les cellules amacrines en étoile (\nt{GABA}) les inhibent de façon asymétrique, seulement par les prolongements orientés du côté \og nul\fg. La formule de $v^*$ est une déduction du chapitre & \cite{Barlow1965,Fried2002} & mesuré \\
6 & L'inhibition shuntante divise la tension au lieu de soustraire~\eqref{eq:shunt} (fig.~\ref{fig:shunt}) & Loi d'Ohm pour un diviseur à trois conductances, avec un potentiel d'équilibre du chlorure proche du repos; pour un neurone sans impulsions & déduction du chapitre & théorie \\
7 & Sur la fréquence des impulsions, le shunt agit par soustraction, et la division vient de l'inhibition associée à un bruit de fond équilibré & Modèle: chez un neurone qui se déclenche, le courant à travers le shunt est presque constant, et l'effet sur la fréquence est soustractif. Expérience: dans des neurones pyramidaux du cortex somatosensoriel du rat (tranches), on a introduit par clamp dynamique des conductances de fond excitatrice et inhibitrice; leur augmentation simultanée et équilibrée divise la pente de la réponse sans décalage notable de la fréquence. Revue: la division par l'activité totale se rencontre dans beaucoup de régions du cerveau & \cite{HoltKoch1997,Chance2002,Carandini2012} & mesuré \\
8 & Pour $\beta>1$, l'inhibition mutuelle choisit un seul gagnant (proposition~\ref{prop:wta}, \eqref{eq:wta}) & Valeurs propres $-(1\pm\beta)/\tau$: déduction du chapitre. Les fréquences $0.73$ et $0.53$ pour $\beta=0.5$ sont le point fixe $r_{1,2}=(I_{1,2}-\beta I_{2,1})/(1-\beta^2)$; pas de script dans \texttt{models/}. Le chapitre ne cite pas d'expérience directe & déduction du chapitre & théorie \\
9 & Effacement de la mémoire par un signal via \nt{GABA}; deux groupes à inhibition mutuelle forment un verrou & Découle de la fig.~\ref{fig:bistable} et de la proposition~\ref{prop:wta}. Pas d'expérience & déduction du chapitre & théorie \\
10 & L'équilibre de la boucle \nt{GLUT}--\nt{GABA} est stable si l'inhibition est assez forte et assez rapide (proposition~\ref{prop:ei}) & Modèle à deux populations~\eqref{eq:ei}; les conditions (i) et (ii) portent sur le déterminant et la trace de sa matrice, déduits dans le chapitre & \cite{WilsonCowan1972} & théorie \\
11 & Pour $w_{EE}=1.5$ et $w_{EI}w_{IE}=2$, une inhibition rapide ($\tau_I=2\ms$) maintient $E^*=0.67h$, une inhibition lente ($30\ms$) déclenche des oscillations (fig.~\ref{fig:ei}) & Résolution numérique de~\eqref{eq:ei}, script \texttt{figures/make\_ei.py} & calcul & modèle \\
12 & Le cortex fonctionne en régime stabilisé par l'inhibition; signature: poussez les neurones \nt{GABA}, et leur fréquence baisse~\eqref{eq:isn} & La théorie a prédit la réponse paradoxale. Expérience: enregistrements intracellulaires dans le cortex visuel primaire du chat; quand la réponse est supprimée par un stimulus environnant, la conductance inhibitrice n'augmente pas, mais baisse avec la conductance excitatrice. L'excitation optogénétique des neurones PV de la souris dans les cortex visuel, somatosensoriel et moteur abaisse la fréquence des neurones inhibiteurs, y compris sans stimulus et sous anesthésie légère. Le coefficient $-1/3$ avec les valeurs de la fig.~\ref{fig:ei} est une déduction du chapitre & \cite{Tsodyks1997isn,Ozeki2009,Sanzeni2020} & mesuré \\
13 & La boucle \og \nt{GLUT} excite \nt{GABA}, \nt{GABA} inhibe \nt{GLUT}\fg{} engendre un rythme rapide de période $12$--$50\ms$ & La stimulation optogénétique des neurones \nt{GABA} rapides dans le cortex somatosensoriel de la souris, à des fréquences de $8$--$200$\,Hz, renforce sélectivement le rythme gamma ($20$--$80$\,Hz, période $12$--$50\ms$); la stimulation des neurones \nt{GLUT} ne renforce que des rythmes plus lents & \cite{Cardin2009,Buzsaki2012} & mesuré \\
\bottomrule
\end{longtable}}
